\subsection{Environmental context constrains historical attribution}

The fossil measurements describe cranial volume, rather than a direct measure of cognition or partner assessment. Gingerich's consensus workbook supplies \MwcFossilN{} observations and identifies approximately \MwcStasisOlder--\MwcStasisYounger{} million years ago (Ma) as a published interval of relative stasis \citep{gingerich2022}. Figure~\ref{fig:minds-environment} shows observations and 100-kyr geometric means without imposing a flat trajectory. Apparent phases depend on temporal scale, sampling and model specification. Du and colleagues distinguish whole-clade patterns from processes within lineages; the two need not support the same evolutionary description \citep{du2018}. Ecological and demographic conditions therefore constrain attribution of anatomical change to any proposed social mechanism.

The Middle Pleistocene Transition, conventionally approximately \MwcMptOlder--\MwcMptYounger{} Ma, involved a shift toward longer, stronger glacial cycles rather than one precisely dated catastrophe \citep{clark2006}. The published LR04 benthic oxygen-isotope stack provides broader chronological context: its \MwcClimateRecords{} marine records combine information about global ice volume and deep-ocean temperature, not direct air temperature \citep{lisiecki2005}. Marine and terrestrial proxies document pronounced North Atlantic and southwest European cooling during MIS 34, approximately \MwcMisOlder--\MwcMisYounger{} Ma, including an extreme terminal stadial \citep{margari2023}. These observations establish environmental disruption within the later part of the published stasis phase.

Margari and colleagues infer that the cooling likely caused European hominin depopulation, combining proxy evidence, habitat modeling and the occupation record \citep{margari2023}. The demographic consequence is an inference with regional scope: European absence or recolonization does not establish extinction across Africa and Asia. Nor does an occupation gap directly measure the population contributing to later lineages. Environmental change could alter which populations survive, move or enter the fossil sample, complicating an anatomical trajectory even without a uniform change within each lineage.

A separate genetic analysis inferred an ancestral bottleneck at approximately \MwcBottleneckOlder--\MwcBottleneckYounger{} Ma, with effective population size near \MwcEffectiveSize{} \citep{hu2023}. Effective size is a genetic-model parameter, not a census of surviving humans. Deng and colleagues argue that the severe bottleneck can be a statistical artifact; Cousins and Durvasula find insufficient evidence and a better fit from a model without it \citep{deng2025,cousins2025}. Zhou and colleagues' July 2026 defense challenges computational and modeling criticisms, but does not establish consensus \citep{zhou2026}. Muttoni and Kent's distinct March 2024 timing-and-migration paper concerns MIS 22 near 0.9 Ma; it is neither the 2026 defense nor the MIS 34 cooling study \citep{muttoni2024}.

Australasian tektites establish an impact near \MwcImpactAge{} Ma \citep{sieh2023}. Arguments locating its source under the Bolaven volcanic field require separate geological assessment; published criticism questions that location \citep{mizera2022}. No verified evidence here establishes that the impact caused cranial-volume stasis. The impact, cooling intervals and proposed bottleneck all postdate the published phase's onset. They may contextualize its later portion or subsequent transitions, but overlap cannot explain the entire interval or identify a cause.

Mechanisms also make competing predictions. Resource stress could constrain costly development and favor energy-saving responses; environmental challenges could favor behavioral flexibility without requiring larger cranial volume. Migration, extinction and changing sampled lineages could shift the aggregate trajectory. Will and colleagues found a strong temperature association with body size, while brain-size associations were smaller and involved other environmental variables, limiting a simple cooling-to-larger-brain claim \citep{will2021}. Distinguishing these hypotheses requires geographically matched chronology, developmental and resource evidence, and within-lineage comparisons robust to sampling and dating assumptions. Establishing partner assessment's contribution would additionally require evidence for its own consequential links. The present alignment estimates none of those contributions.

\mwcfigure{minds-environment.pdf}{fig:minds-environment}{\textbf{Brain-volume trajectories and environmental context.} A: \MwcFossilZoomN{} of \MwcFossilN{} consensus observations in the enlarged window; gray points are specimens and black symbols are 100-kyr geometric means. Empty bins remain gaps. Shading marks Gingerich's published phase, not a fitted plateau; the inset shows the complete trajectory. The workbook supplies no specimen dating errors. B: the LR04 stack on its published sampling and age model; benthic $\delta^{18}$O (\textperthousand) reflects ice volume and deep-ocean temperature \citep{lisiecki2005}. C: solid bands show approximate climate-interval durations, the diamond an impact date, and the dashed hatched bar a disputed, model-inferred demographic duration. Neither band endpoints nor the bottleneck duration are confidence intervals; analytical dating uncertainties are tabulated separately in the supplement. MIS 34 cooling is proxy evidence; European hominin depopulation is Margari and colleagues' regional inference \citep{margari2023}. Alignment supplies chronological context rather than a causal estimate.}{environment}
